\section{Construction of the recursive radial frame}
\label{re-section}

The terminal commutation and determinant identities suffice to
construct the change of variables in the positivity argument.
We first construct its initial connection directly from a
polynomial coefficient.  We then prove the two recursive changes
of variables by \(2\times2\) matrix multiplication.  In particular,
the larger normal form involving individual raising operators is
not a premise of this construction.

\subsection{The initial connection}

Let \(k\geq1\), let \(W\) have signature \((k,1)\), and retain
the negative vectors \(z,z'\), with
\[
 h(z,z)=-T,\qquad h(z',z')=-T',\qquad
 \zeta(m)=h(mz',z),\qquad G=\zeta\bar\zeta-TT'.
\]
We work on the open set \(G>0\).  In particular, \(G,\zeta,\bar\zeta\)
are invertible smooth multiplication operators.  The space of
smooth functions has the fixed left and right compact equivariance
of Section~\ref{rd-section}; it can also depend on independent
variables \(c_1,\ldots,c_r\in\mathbb C^\times\).
The formal adjoint is taken for terminal Haar measure, the
Hilbert--Schmidt inner product, and
\(\prod_i |c_i|^{-2}\,dc_i\).

Write
\[
 t(Y)=
 \begin{pmatrix} a(Y)&b(Y)\\c(Y)&\tau(Y)\end{pmatrix}
\]
for the actual terminal matrix of
\eqref{rc-t21}--\eqref{rc-t22}, including its full Capelli
entry \(a(Y)\).  The letter \(c(Y)\) in this display denotes
an operator, and is unrelated to the independent coordinates
\(c_i\).  Sections~\ref{rq-section} and \ref{rr-section} give
\begin{align}
 (u-v)[t_{ij}(u),t_{kl}(v)]
 &=t_{kj}(u)t_{il}(v)-t_{kj}(v)t_{il}(u),       \label{re-rtt}\\
 a(Y+1)\tau(Y)-c(Y+1)b(Y)&=q(Y),                \label{re-one-det}\\
 q(Y)&=-TT'\prod_{i=1}^k
       (Y+\tfrac12+u_i)(Y+\tfrac12+u'_i).       \label{re-q}
\end{align}
The real numbers \(u_i,u'_i\) are the shifted compact highest
weights used in the terminal determinant calculation.
All these equalities are equalities of differential operators,
not merely of their principal symbols.

\begin{lemma}\label{re-base-leading}
Put \(h=k-1\).  There are operators \(B,C\) such that
\begin{align}
 \tau(Y)&=Gx_0(Y),&
 x_0(Y)&=Y^h+B Y^{h-1}+O(Y^{h-2}),             \label{re-x-leading}\\
 c(Y)&=-\bar\zeta Y^{h+1}+C Y^h+O(Y^{h-1}).    \label{re-c-leading}
\end{align}
Here \(x_0(Y)=\sum_\rho p_\rho(Y)P_\rho\), with
\(p_\rho(Y)=\prod_{i=1}^{k-1}(Y+\rho_i)\), is the actual
finite spectral decomposition of
Theorem~\ref{rd-constructed-kappa}.
When \(h=0\), the first formula means \(x_0=1\), we set
\(B=0\), and the second formula has no remainder.
The operator \(C\) has differential order at most one.
For real \(Y\), the formal adjoint of \(t_{ij}(Y)\) is
\(t_{ji}(Y)\).
\end{lemma}

\begin{proof}
The spectral formula, including the fact that its coefficients
have order zero, was proved in
Theorem~\ref{rd-constructed-kappa}.
There is also a direct check on its leading normalization.
Put \(y=m^{-1}z\), \(\eta_i=h(mw_i,z)\),
\(x=z'\), and \(\xi_i=h(w_i,z')\).
The leading coefficient of the double cofactor contraction is
\[
 -\sum_{i<j}(\eta_i\xi_j-\eta_j\xi_i)(y_i x_j-y_jx_i)
 =-\{(\eta y)(\xi x)-(\eta x)(\xi y)\}=G.
\]
The leading coefficient of the single cofactor contraction is
\(-\xi y=-\bar\zeta\).  A cofactor of size \(k\) has degree
\(k\) in its parameter; its coefficient of degree \(k-1\)
contains at most one Lie algebra generator.  The shift
\(Y\mapsto Y+k/2\) changes this coefficient only by an order
zero term.  This proves \eqref{re-c-leading} and the order
assertion.

For completeness, the adjoint assertion follows directly from
the same minors.  In an orthonormal basis put
\(\epsilon_i=h(w_i,w_i)\in\{1,-1\}\).
Complexified right differentiation satisfies
\(E_{ij}^*=\epsilon_i\epsilon_jE_{ji}\).
The row and column forms of a quantum minor give
\[
 Q_{I,J}(u)^*=\epsilon_I\epsilon_J Q_{J,I}(\bar u).
\]
Indeed, taking adjoints reverses the order of all factors;
use the row form with both index lists reversed.  Reversing
both lists contributes the square of the same permutation sign,
which is one.  For the signed cofactors this reads
\(C_{ij}(u)^*=\epsilon_i\epsilon_jC_{ji}(\bar u)\).
Consequently, for real \(u\),
\[
 \left(-\sum_{i,j}\xi_i y_j C_{ij}(u)\right)^*
 =-\sum_{i,j}C_{ij}(u)\eta_i x_j
 =-\sum_{i,j}\eta_i x_j C_{ij}(u-1).
\]
The last equality is precisely the covector shift
\eqref{rq-shift-one}.  Thus \(c(Y)^*=b(Y)\).
The full determinant is self-adjoint by the same minor rule.
Finally \(\tau(Y)=Gx_0(Y)\) is self-adjoint on the radial
space: \(G\) is real and each \(P_\rho\) is a pointwise
orthogonal projection.  This proves every matrix entry of the
assertion.
\end{proof}

\begin{lemma}\label{re-two-dets}
The relations \eqref{re-rtt} and \eqref{re-one-det} imply
\begin{equation}
 t(Y+1)^t Jt(Y)=q(Y)J,\qquad
 t(Y)Jt(Y+1)^t=q(Y)J,\qquad
 J=\begin{pmatrix}0&1\\-1&0\end{pmatrix}.       \label{re-dets}
\end{equation}
They also imply that the members of each fixed entry family
\(t_{ij}(Y)\) commute with one another.
\end{lemma}

\begin{proof}
For a fixed entry, \eqref{re-rtt} gives
\((u-v-1)[t_{ij}(u),t_{ij}(v)]=0\).
Interchanging \(u,v\) gives the same assertion with
\(u-v+1\); subtraction proves commutation, without a restriction
on \(u-v\).
At \(u=Y+1,v=Y\), the remaining entries of
\eqref{re-rtt} give
\[
 a_+c=c_+a,\quad b_+\tau=\tau_+b,\quad
 ab_+=ba_+,\quad c\tau_+=\tau c_+,
\]
where a subscript \(+\) means evaluation at \(Y+1\).
They further give
\[
 a_+\tau-c_+b=a\tau_+-bc_+
                         =\tau a_+-cb_+
                         =\tau_+a-b_+c.
\]
For example, the first equality is the relation for \(c_+,b\);
the second follows from that for \(a_+,\tau\).
The fourth expression follows in the same way from that for
\(b_+,c\).  Multiplying the two displayed \(2\times2\)
matrices now gives \eqref{re-dets} entry by entry.
\end{proof}

\begin{theorem}\label{re-initial-connection}
Define
\begin{equation}
 \nabla=-G^{-1}C-(B+\tfrac12)\bar\zeta G^{-1},
 \qquad
 \nabla^*=\text{its formal adjoint}.
                                                               \label{re-connection}
\end{equation}
Then \(\nabla\) is a first-order operator and
\begin{gather}
 [\nabla,\zeta]=1,\quad[\nabla,\bar\zeta]=0,\quad
 [\nabla,G]=\bar\zeta,\qquad
 [x_0(Y),\nabla]=0,                               \label{re-connection-laws}\\
 [\nabla^*,\bar\zeta]=-1,\quad[\nabla^*,\zeta]=0,\quad
 [\nabla^*,G]=-\zeta.                            \label{re-star-laws}
\end{gather}
The connection and its adjoint commute with the independent
coordinates \(c_i,\bar c_i\), their inverses, and their Euler
derivatives.

Let \(\sigma\in\{1,-1\}\), replace \(t\) by \(\sigma t\),
and put \(x(Y)=\sigma x_0(Y)\).  Define
\begin{equation}
 I=\{Y\in\mathbb C:p_\rho(Y)\ne0\text{ for every }\rho\},\quad
 x(Y)^{-1}=\sigma\sum_\rho p_\rho(Y)^{-1}P_\rho\quad(Y\in I).
                                                               \label{re-initial-domain}
\end{equation}
On \(I\) put
\begin{equation}
 \ell(Y)=(Gx(Y))^{-1}(\sigma c(Y)),\qquad
 \ell_0(Y)=-\nabla-(Y+\tfrac12)\bar\zeta G^{-1},\qquad
 e(Y)=\zeta G\bar\zeta^{-1}\{\ell(Y)-\ell_0(Y)\}.
                                                               \label{re-initial-e}
\end{equation}
Each \(e(Y)\) commutes with \(\zeta,\bar\zeta,G\) and their
inverses.  These operators supply the actual initial data of
the recursive construction, with
\begin{equation}
 \begin{aligned}
 &R=F=1,\quad N=0,\quad d=c_{\rm wt}=1,\quad t_{22}(Y)=Gx(Y),\\
 &t_{21}(Y)=Gx(Y)\{\ell_0(Y)+\bar\zeta\zeta^{-1}G^{-1}e(Y)\}.
 \end{aligned}
                                                               \label{re-initial-frame}
\end{equation}
The determinant \(q\) is unchanged by the sign \(\sigma\).
\end{theorem}

\begin{proof}
Compare coefficients of the polynomial parameter \(v\) in
\eqref{re-rtt}.  Commutation of the family \(c(v)\) gives
\([c(u),\bar\zeta]=0\).
The coefficient of \(v^{h+2}\) in the relation for
\(\tau(u),c(v)\) gives \([\tau(u),\bar\zeta]=0\).
Commutation of the family \(\tau(v)\) gives
\([\tau(u),G]=0\).  Since \(G=\zeta\bar\zeta-TT'\) and
\(\bar\zeta\) is invertible, also \([\tau(u),\zeta]=0\).
Finally the coefficient of \(v^{h+1}\) in
\[
 (u-v)[c(u),\tau(v)]
                    =c(u)\tau(v)-c(v)\tau(u)
\]
is
\([c(u),G]=-\bar\zeta\tau(u)\).
It follows that
\begin{equation}
 [c(u),\zeta]=-\tau(u),\qquad [c(u),\bar\zeta]=0.
                                                               \label{re-zeta-comm}
\end{equation}
Comparing the coefficient of \(u^h\) proves
\([C,\zeta]=-G\), \([C,\bar\zeta]=0\).  The coefficient \(B\)
commutes with scalar multiplication by the spectral
decomposition.  Thus \eqref{re-connection} gives the first
three identities in \eqref{re-connection-laws}.

It remains to check the commutation with \(x_0(u)\).
Work in the algebra of formal Laurent series in \(v^{-1}\)
with differential-operator coefficients.  The leading term
of \(\tau(v)\) is \(Gv^h\), with invertible coefficient, so
its formal inverse exists.  Equations
\eqref{re-x-leading}--\eqref{re-connection} give
\begin{equation}
 \ell(v)=\tau(v)^{-1}c(v)
      =-v\bar\zeta/G-\nabla-\tfrac12\bar\zeta/G+O(v^{-1}).
                                                               \label{re-laurent}
\end{equation}
Dividing the RTT relation for \(\tau(u),c(v)\) by the
commuting \(\tau(v)\) gives
\begin{equation}
 [\tau(u),\ell(v)]
       =-\tau(u)\frac{\ell(u)-\ell(v)}{u-v}.       \label{re-tau-ell}
\end{equation}
This identity may initially be read with \(\tau(u)\) also
formally invertible; multiplication by it cancels its inverse.
The coefficient of \(v^0\) on its right is
\(\bar\zeta x_0(u)\).  On its left the coefficient is
\[
 [Gx_0(u),-\nabla]
        =-G[x_0(u),\nabla]+\bar\zeta x_0(u).
 \]
Hence \([x_0(u),\nabla]=0\).  This is a polynomial
identity in \(u\); it therefore holds for every complex \(u\).
The interpolation formula for \(P_\rho\) then also gives
\([P_\rho,\nabla]=0\) for every \(\rho\).
Taking formal adjoints gives \eqref{re-star-laws}.
All coefficients used here act only in the terminal variable,
which proves the assertions about the independent \(c_i\).

The finite spectral formula proves the stated two-sided
inverse on \(I\), commutation of the \(x(Y)\), their
commutation with all scalar functions, and their
self-adjointness for real \(Y\).
Equation \eqref{re-zeta-comm} shows
\([\ell(Y),\zeta]=-1\) and
\([\ell(Y),\bar\zeta]=0\).
The same is true of \(\ell_0(Y)\); their difference therefore
commutes with \(\zeta,\bar\zeta,G\), proving the assertion about
\(e\).  The last displayed identities are now definitions.
The sign disappears from \(\ell\), and its square in the
determinant is one.  For \(k=1\), \eqref{re-laurent} is an
exact equality, so \(e=0\).  No limiting argument at a spectral
root has been used.
\end{proof}

The construction stays within the actual radial function space.
The cofactor covariance in Section~\ref{rc-section} shows that
the terminal entries preserve the two compact equivariances.
Each of their polynomial coefficients is a finite linear
combination of evaluations, so \(C\) preserves them as well.
The spectral projections preserve them by
Theorem~\ref{rd-constructed-kappa}; multiplication by
\(\zeta,\bar\zeta,G\) and their inverses does so because these
functions are invariant.  All these coefficients are smooth on
\(G>0\).  Formal adjoints preserve the equivariance because
the compact actions are unitary.  Consequently every operator
just constructed acts on the stated smooth radial space and
preserves supports.

\subsection{Two elementary Gauss identities}

The following calculations take place in an arbitrary
associative complex algebra.  Every inverse mentioned is a
two-sided inverse.  Suppose \(t(u)\) satisfies
\eqref{re-rtt} and \eqref{re-dets}, and that \(\tau(u)\) is
invertible for \(u\) in a set \(I\).  Write
\[
 \ell_u=\tau_u^{-1}c_u,\qquad
 m_u=\tau_u^{-1}b_u,\qquad
 H_u=q(u)\tau_u^{-1}\tau_{u+1}^{-1}.
\]
The last expression is used only when \(u,u+1\in I\).

\begin{lemma}\label{re-gauss}
Whenever the indicated adjacent points belong to \(I\),
\begin{align}
 m_{u+1}&=b_u\tau_u^{-1},&
 \ell_u&=c_{u+1}\tau_{u+1}^{-1},               \label{re-shifts}\\
 a_u-b_u\tau_u^{-1}c_u&=q(u)\tau_{u+1}^{-1},&
 a_u-c_u\tau_u^{-1}b_u&=q(u-1)\tau_{u-1}^{-1},\label{re-schur}\\
 t(u)&=
 \begin{pmatrix}1&m_{u+1}\\0&1\end{pmatrix}
 \begin{pmatrix}q(u)\tau_{u+1}^{-1}&0\\0&\tau_u\end{pmatrix}
 \begin{pmatrix}1&0\\\ell_u&1\end{pmatrix},     \label{re-gauss-plus}\\
 t(u)^t&=
 \begin{pmatrix}1&\ell_{u-1}\\0&1\end{pmatrix}
 \begin{pmatrix}q(u-1)\tau_{u-1}^{-1}&0\\0&\tau_u\end{pmatrix}
 \begin{pmatrix}1&0\\m_u&1\end{pmatrix}.        \label{re-gauss-minus}
\end{align}
If \(u\ne v\), put
\(\eta_{uv}=(\ell_u-\ell_v)/(u-v)\) and
\(\mu_{uv}=(m_u-m_v)/(u-v)\).  Then
\begin{align}
 [\tau_u,\ell_v]&=-\tau_u\eta_{uv},&
 [\tau_u,m_v]&=\tau_u\mu_{uv},                 \label{re-gauss-comm}\\
 [\ell_u,\ell_v]&=\eta_{uv}(\ell_u-\ell_v),    \label{re-ell-comm}\\
 (u-v-1)[m_u,\ell_v]&=H_v-H_{u-1}.            \label{re-mixed}
\end{align}
For \eqref{re-mixed} assume \(u-1,u,v,v+1\in I\).
If \(d=d^*\) is invertible and \(dt(u)\) is Hermitian for real
\(u\), then
\begin{equation}
             dm_{u+1}=\ell_u^*d                 \label{re-adj-gauss}
\end{equation}
for real \(u,u+1\in I\).
\end{lemma}

\begin{proof}
The zero diagonal entries of the two identities
\eqref{re-dets} give
\(b_{u+1}\tau_u=\tau_{u+1}b_u\) and
\(c_u\tau_{u+1}=\tau_uc_{u+1}\).
These are \eqref{re-shifts}.  The off-diagonal entries,
with these substitutions, give \eqref{re-schur}.
Multiplying the three factors proves the two factorizations.

The RTT relation for \(\tau_u,c_v\), after multiplication
by \(\tau_v^{-1}\), gives the first identity in
\eqref{re-gauss-comm}.  For clarity, the second can be
derived without dividing by \(u-v+1\).
Put \(D=u-v\) and \(K=[\tau_u,m_v]\).
The relations for \(\tau_u,b_v\), together with the same
relation with \(u,v\) interchanged, give
\[
 DK=\tau_v^{-1}b_u\tau_v-m_v\tau_u,\qquad
 \tau_v^{-1}b_u\tau_v=b_u-K.
\]
Thus \((D+1)K=b_u-m_v\tau_u
=b_u-\tau_um_v+K\), proving \(DK=\tau_u(m_u-m_v)\).
Next expand \(c_uc_v=c_vc_u\), use
\(c_w=\tau_w\ell_w\), and move \(\tau_u,\tau_v\)
to the left using \eqref{re-gauss-comm}.
Cancellation of their invertible product gives
\eqref{re-ell-comm}.

The relation for \(b_u,c_v\) gives, by the same movement,
\[
 D[m_u,\ell_v]-(m_u-m_v)\ell_v-(\ell_u-\ell_v)m_u
       =H_v-H_{u-1}+m_v\ell_v-\ell_um_u.
\]
Here use the first Schur expression at \(v\) and the second
at \(u\), so that no inverse at \(u+1\) is required.
The left side is
\((D-1)[m_u,\ell_v]+m_v\ell_v-\ell_um_u\).
Cancellation proves \eqref{re-mixed}.
Finally weighted Hermiticity says \(c_u^*d=db_u\) and
\(\tau_u^*d=d\tau_u\).  Hence
\[
 \ell_u^*d=c_u^*(\tau_u^*)^{-1}d
                    =db_u\tau_u^{-1}=dm_{u+1}.
\]
\end{proof}

\subsection{The recursive construction and its domains}

We now allow an auxiliary tensor product of finitely many copies
of \(\mathbb C^2\).  All matrices below have differential-operator
coefficients on the radial space and act on this tensor product.
A new factor always means a factor on which every existing
coefficient acts as the identity.  This condition is essential.
Let
\[
 w=\zeta\bar\zeta/G,\qquad
 \kappa=\bar\zeta/(\zeta G),\qquad t_z=TT'/\bar\zeta.
\]
A frame consists of sets \(I\subset\Omega\), a matrix \(t(Y)\)
defined on \(\Omega\), an invertible \(F\), commuting operators
\(x(Y)\), an invertible self-adjoint \(c_{\rm wt}\), and
operators \(e(Y)\) for \(Y\in I\).  Its scalar determinant
\(q(Y)\) is defined on the adjacent-point domain
\(\Omega^{(2)}=\{Y:Y,Y+1\in\Omega\}\).  Put
\begin{equation}
 d=w^N c_{\rm wt},\quad
 \ell_Y=-\nabla-(Y+\tfrac12)\bar\zeta/G+\kappa e(Y),\quad
 m_Y=d^{-1}\ell_{Y-1}^*d
       \quad(Y\in\mathbb R,\ Y-1\in I).
                                                               \label{re-frame-data}
\end{equation}
Here \(\ell_Y\) is defined for \(Y\in I\).  For real \(Y\)
with \(Y-1,Y\in I\), the displayed \(m_Y\) agrees with
the upper Gauss coordinate by the adjoint identity
\eqref{re-adj-gauss}.
The frame identities are
\begin{align}
 R&=F^*dF,&
 Ft_{22}(Y)F^{-1}&=Gx(Y)\quad(Y\in\Omega),       \label{re-frame-main}\\
 Ft_{21}(Y)F^{-1}&=Gx(Y)\ell_Y
                    \quad(Y,Y+1\in I).        \label{re-frame-lower}
\end{align}
We require \(x(Y)^{-1}\) on \(I\), that the \(x(Y)\)
commute with one another, \(c_{\rm wt}\), \(\nabla\), and
\(\zeta,\bar\zeta,G\), that \(x(Y)^*=x(Y)\) for real \(Y\),
and that \(c_{\rm wt}\) and \(e(Y)\) commute with the latter
three multiplication operators.  In addition, \(R=R^*\),
\(Rt(Y)\) is Hermitian at real \(Y\in\Omega\), and
\eqref{re-rtt} holds on \(\Omega\), with
\eqref{re-dets} whenever \(Y,Y+1\in\Omega\).
These conditions hold for the initial frame constructed above.

For a new factor let
\begin{equation}
 P_{ij}=E_{ji},\qquad Q_{ij}=E_{ij},\qquad
 R_P(D)=\frac{DI-P}{D-1},\qquad
 R_Q(d_0)=\frac{d_0I-Q}{d_0-1}.
                                                               \label{re-fresh-matrices}
\end{equation}
Here \(i,j\in\{1,2\}\) are the matrix indices of \(t\),
and \(E_{ij}\) acts on the new factor.  The operation
\(\operatorname{ins}t(a)\) means that the two indices of
\(t(a)\) act on that new factor.

\begin{theorem}\label{re-recursive-frame}
Suppose the preceding frame conditions hold.

For a real number \(a\), assume
\begin{equation}
                a,a+1\in I,\qquad q(a)\ne0.     \label{re-plus-hyp}
\end{equation}
There is a new frame with
\begin{equation}
 \begin{aligned}
 \Omega_+&=\Omega\setminus\{a+1\},&
 I_+&=I\setminus\{a,a+1\},\\
 t_+(Y)&=t(Y)R_P(Y-a),&
 q_+(Y)&=q(Y)\frac{Y-a+1}{Y-a},\\
 R_+&=R\,\operatorname{ins}t(a),&
 N_+&=N+1.
 \end{aligned}                                                   \label{re-plus-module}
\end{equation}
Its other entries are
\begin{equation}
 F_+=C_+L_aF,\quad
 C_+=\begin{pmatrix}\zeta^{-1}&0\\0&G/\bar\zeta\end{pmatrix},\quad
 L_a=\begin{pmatrix}1&0\\\ell_a&1\end{pmatrix}.                 \label{re-plus-F}
\end{equation}
\begin{align}
 x_+(Y)&=x(Y)
       \begin{pmatrix}(Y-a)/(Y-a-1)&0\\0&1\end{pmatrix},\nonumber\\
 c_{{\rm wt},+}&=c_{\rm wt}
       \begin{pmatrix}q(a)x(a+1)^{-1}&0\\0&x(a)\end{pmatrix},
                                                               \label{re-plus-xc}
\end{align}
With \(\Delta_{Ya}e=(e(Y)-e(a))/(Y-a)\), put
\begin{equation}
 e_+(Y)=
 \begin{pmatrix}
 e(Y)-\Delta_{Ya}e+t_z&-(Y-a)^{-1}\\
 0&x(a)^{-1}e(Y)x(a)
 \end{pmatrix}.                                \label{re-plus-e}
\end{equation}

For a real number \(b\), assume
\begin{equation}
             b-1,b,b+1\in I,\qquad q(b-1)\ne0.  \label{re-minus-hyp}
\end{equation}
There is a new frame with
\begin{equation}
 \begin{aligned}
 \Omega_-&=\Omega\setminus\{b-1\},&
 I_-&=I\setminus\{b-1,b\},\\
 t_-(Y)&=R_Q(b-Y)t(Y),&
 q_-(Y)&=q(Y)\frac{b-Y}{b-Y-1},\\
 R_-&=R\,\operatorname{ins}t(b)^t,&
 N_-&=N+1.
 \end{aligned}                                                   \label{re-minus-module}
\end{equation}
Its other entries are
\begin{equation}
 F_-=C_-M_bF,\quad
 C_-=\begin{pmatrix}\bar\zeta^{-1}&0\\0&G/\zeta\end{pmatrix},\quad
 M_b=\begin{pmatrix}1&0\\m_b&1\end{pmatrix}.                  \label{re-minus-F}
\end{equation}
\begin{align}
 x_-(Y)&=x(Y)
       \begin{pmatrix}(b-Y)/(b-Y-1)&0\\0&1\end{pmatrix},\nonumber\\
 c_{{\rm wt},-}&=c_{\rm wt}
       \begin{pmatrix}q(b-1)x(b-1)^{-1}&0\\0&x(b)\end{pmatrix},
                                                               \label{re-minus-xc}
\end{align}
\begin{equation}
 e_-(Y)=
 \begin{pmatrix}
 e(Y)&0\\
 -\dfrac{q(b-1)x(b)^{-1}x(b-1)^{-1}}{b-Y-1}
       &x(b)^{-1}e(Y)x(b)-\Delta_{Yb}e+t_z
 \end{pmatrix}.                                  \label{re-minus-e}
\end{equation}
All displayed rational formulas are asserted on their indicated
domains.  In particular, \(q_\pm\) is defined on
\(\Omega_\pm^{(2)}=\{Y:Y,Y+1\in\Omega_\pm\}\).
For \(q_+\), this condition excludes \(Y=a\), so its displayed
formula never evaluates its apparent pole.
\end{theorem}

\begin{proof}
Conjugate the old matrix by \(F\), and abbreviate it by \(t\)
during the calculation.  Its lower-right entry is
\(\tau_Y=Gx(Y)\); the old density is \(d\).
Similarity preserves \eqref{re-rtt} and \eqref{re-dets}.
The Gauss coordinates coincide with \(\ell_Y\) when
\(Y,Y+1\in I\), and \eqref{re-adj-gauss} identifies the upper
coordinate with \(m_Y\) at the real parameters used below.

First we check the conjugation of the last row.  For the plus
step put \(D=Y-a\) and
\(\eta=(\ell_Y-\ell_a)/D\).  Before conjugating by \(L_a\),
the new entries are
\begin{equation}
 (t_+)_{22}=\frac{\tau_Y}{D-1}
      \begin{pmatrix}D&0\\-\ell_Y&D-1\end{pmatrix},\qquad
 (t_+)_{21}=\frac{\tau_Y}{D-1}
      \begin{pmatrix}(D-1)\ell_Y&-1\\0&D\ell_Y\end{pmatrix}.
                                                               \label{re-plus-blocks}
\end{equation}
Set the two matrices on the right equal to \(U,V\).
The first identity of \eqref{re-gauss-comm} gives
\begin{equation}
 L_a\tau_Y U
       =\tau_Y\operatorname{diag}(D,D-1)L_a.
                                                               \label{re-plus-tau-calc}
\end{equation}
Moreover
\begin{equation}
 V L_a^{-1}
 =U L_a^{-1}
       \begin{pmatrix}\ell_Y-\eta&-D^{-1}\\0&\ell_Y+\eta\end{pmatrix}.
                                                               \label{re-plus-low-calc}
\end{equation}
Three entries follow by \(D\eta=\ell_Y-\ell_a\).
For the remaining lower-left entry the difference is a scalar
multiple of
\((\ell_Y-\ell_a)^2-D[\ell_Y,\ell_a]\), which is zero
by \eqref{re-ell-comm}.
Thus conjugation by \(L_a\) gives the asserted new
\(\tau\) and the Gauss coordinate in
\eqref{re-plus-low-calc}.

The assertion about \(\tau\) holds on all of \(\Omega_+\),
including points where \(x(Y)\) is not invertible.
To see this without division by \(D\), use the RTT identity
in the form
\begin{equation}
 D\ell_a\tau_Y-c_Y=(D-1)\tau_Y\ell_a .
                                                               \label{re-plus-all-reg}
\end{equation}
It requires only the inverse of \(\tau_a\), and proves
\eqref{re-plus-tau-calc} directly with \(c_Y\) in place of
\(\tau_Y\ell_Y\).

For the minus step put \(d_0=b-Y\).
Before conjugating by \(M_b\), the two new entries, multiplied
by \(d_0-1\), are
\[
 A=\begin{pmatrix}d_0\tau_Y&0\\-b_Y&(d_0-1)\tau_Y\end{pmatrix},
 \qquad
 B_0=\begin{pmatrix}d_0c_Y&0\\-a_Y&(d_0-1)c_Y\end{pmatrix}.
\]
The first Schur formula gives
\begin{equation}
 B_0=A
       \begin{pmatrix}\ell_Y&0\\-H_Y/(d_0-1)&\ell_Y\end{pmatrix}.
                                                               \label{re-minus-block}
\end{equation}
The second identity of \eqref{re-gauss-comm} gives
\[
 M_b A=\operatorname{diag}(d_0\tau_Y,(d_0-1)\tau_Y)M_b.
\]
The mixed identity \eqref{re-mixed}, with \(u=b,v=Y\), then gives
\begin{equation}
 M_b
 \begin{pmatrix}\ell_Y&0\\-H_Y/(d_0-1)&\ell_Y\end{pmatrix}
 =
 \begin{pmatrix}\ell_Y&0\\-H_{b-1}/(d_0-1)&\ell_Y\end{pmatrix}M_b.
                                                               \label{re-minus-low-calc}
\end{equation}
For the lower-right entry on all of \(\Omega_-\), use instead
the undivided RTT equality
\begin{equation}
 d_0m_b\tau_Y-b_Y=(d_0-1)\tau_Ym_b.             \label{re-minus-all-reg}
\end{equation}
It only uses \(\tau_b^{-1}\).

It remains to conjugate by \(C_+\) or \(C_-\).
The level relations give
\begin{equation}
 \eta=-\bar\zeta/G+\kappa\Delta_{Ya}e,\qquad
 \bar\zeta/G-\zeta^{-1}=\kappa t_z .
                                                               \label{re-coordinate-simplify}
\end{equation}
The first relation in \eqref{re-gauss-comm}, applied to
\(\tau_a=Gx(a)\), also gives
\begin{equation}
               x(a)^{-1}e(Y)x(a)=e(Y)+\Delta_{Ya}e.
                                                               \label{re-e-conj}
\end{equation}
Here \(Y,Y+1,a,a+1\in I\) and \(Y\ne a\), exactly as in the
new lower-entry domain.
For the plus step, the two diagonal entries after conjugation
are
\[
 \ell_Y-\eta-\zeta^{-1},\qquad
 \ell_Y+\eta+\bar\zeta/G,
\]
and the upper entry is \(-\kappa/D\).
Equations \eqref{re-coordinate-simplify} and
\eqref{re-e-conj} give exactly \eqref{re-plus-e}.
For the minus step the diagonal entries are
\(\ell_Y\) and \(\ell_Y+\bar\zeta/G-\zeta^{-1}\).
Its lower entry is
\[
 -\frac{G}{\zeta}\frac{H_{b-1}}{d_0-1}\bar\zeta
 =-\kappa\,\frac{q(b-1)x(b)^{-1}x(b-1)^{-1}}{d_0-1}.
\]
Using \eqref{re-e-conj} with \(a=b\) gives
\eqref{re-minus-e}.  The requirement \(b+1\in I\) is used
in this last application.

We next prove the congruence for the form.  The weighted
Gauss factorization gives
\begin{align*}
 d\,t(a)
 &=L_a^*\operatorname{diag}
       \bigl(dq(a)\tau_{a+1}^{-1},d\tau_a\bigr)L_a,\\
 d\,t(b)^t
 &=M_b^*\operatorname{diag}
       \bigl(dq(b-1)\tau_{b-1}^{-1},d\tau_b\bigr)M_b.
\end{align*}
All \(x\)'s commute with \(d\) and with the coordinate
multiplications.  Since
\(w/(\zeta\bar\zeta)=G^{-1}\), direct diagonal multiplication
gives
\begin{align*}
 C_+^*w^{N+1}c_{{\rm wt},+}C_+
 &=\operatorname{diag}(dq(a)\tau_{a+1}^{-1},d\tau_a),\\
 C_-^*w^{N+1}c_{{\rm wt},-}C_-
 &=\operatorname{diag}(dq(b-1)\tau_{b-1}^{-1},d\tau_b).
\end{align*}
Undoing the old conjugation therefore proves
\(R_\pm=F_\pm^*w^{N+1}c_{{\rm wt},\pm}F_\pm\).
The inverse of \(F_+\), for example, is
\(F^{-1}L_a^{-1}C_+^{-1}\); the minus inverse is identical
with \(M_b,C_-\) substituted.

We verify the remaining matrix relations explicitly.
The matrices in \eqref{re-fresh-matrices} satisfy
\eqref{re-rtt}: their numerators have, respectively, the
generator matrices \(-E_{ji}\) and \(E_{ij}\), and
\([E_{ij},E_{kl}]=\delta_{jk}E_{il}-\delta_{il}E_{kj}\)
gives each of the sixteen identities.
Multiplication gives their two shifted determinants as
\((D+1)/D\) and \(d_0/(d_0-1)\), respectively.
Before division the numerators are
\((D-1)(D+1)\) and \(d_0(d_0-2)\); the successive
denominators are \((D-1)D\) and
\((d_0-1)(d_0-2)\).
Products of matrices with mutually commuting coefficient
algebras preserve RTT: in the two auxiliary tensor copies,
move each old coefficient past each fresh coefficient and
apply the two RTT equations in succession.
They also preserve the two determinant equations, since
\((AB)^t=B^tA^t\) for such coefficient algebras.
This proves the claimed \(q_\pm\).

The new weighted Hermiticity follows from the old RTT
relation at \(a,Y\), or \(b,Y\).  More explicitly, after
using \(dt(u)=(dt(u))^*\), the plus assertion is the exchange
identity
\[
 \operatorname{ins}t(a)\,t(Y)R_P(Y-a)
       =R_P(Y-a)t(Y)\operatorname{ins}t(a);
\]
the minus assertion is
\[
 \operatorname{ins}t(b)^t R_Q(b-Y)t(Y)
       =t(Y)R_Q(b-Y)\operatorname{ins}t(b)^t.
\]
The first is RTT on the old auxiliary and new factors.
For the second, expand \(Q_{ij}=E_{ij}\): its
\((i,j;k,l)\) entry is the same RTT identity with the two
indices belonging to the new factor interchanged.
This is an entry calculation; it does not use a rule
reversing arbitrary products under partial transpose.
The self-adjointness of \(R_\pm\) also follows from the
congruences just proved.  The reality of \(q_\pm\) follows
from the reality of their parameters.
Commutation within the new lower-left family follows from
RTT as in Lemma~\ref{re-two-dets}.

Finally the inverses in the new \(x,c_{\rm wt}\) exist on
the specified sets by \eqref{re-plus-hyp} or
\eqref{re-minus-hyp}.  Their commutations and adjoints
follow entry by entry from the old ones: every new
diagonal factor is a real scalar rational function times
one of the commuting \(x\)'s or their inverses.
The formulas for \(e_\pm\) show that they commute with
the three coordinate multiplications.
All frame conditions have now been checked.
\end{proof}

\subsection{Iteration, scalar densities and the empty word}

Start with the actual frame of
Theorem~\ref{re-initial-connection}, with
\(\sigma=(-1)^{k-l}\), \(1\le l\le k\), as in the application.
Take an ordered sequence of plus and minus steps on distinct
auxiliary factors.  At each step require exactly
\eqref{re-plus-hyp} or \eqref{re-minus-hyp} for the current
domain and current determinant.
Theorem~\ref{re-recursive-frame} then applies by induction.
This formulation allows the sharper recursively admissible
domains of the source.

A convenient sufficient condition, used for its original
admissible words, is the following.  All parameters \(a_j\)
are real, \(a_j-1,a_j,a_j+1\) belong to the initial set \(I\),
\(q(a_j),q(a_j-1)\ne0\), and
\begin{equation}
             a_i-a_j\notin\{-2,-1,0,1,2\}\quad(i\ne j).
                                                               \label{re-word-domain}
\end{equation}
Indeed an earlier step removes only its parameter and a
neighbor from \(I\).  Its determinant multiplier has its
zero and pole at such neighbors.  Condition
\eqref{re-word-domain} prevents any later required evaluation
from meeting either one.  Distinct auxiliary factors supply
the required commutation with each new factor.
This proves sufficiency for every order of the steps.

\begin{lemma}\label{re-shifted-domain}
For each step define \(y_j=a_j\) in the plus case and
\(y_j=b_j-1\) in the minus case.
It is enough to require that the auxiliary factors are distinct,
that all \(y_j\) are real, that
\[
 y_j,y_j+1\in I,\qquad q(y_j)\ne0,\qquad
 y_j+2\in I\ \hbox{for a minus step},
\]
in the initial frame, and that
\[
 y_i-y_j\notin\{-2,-1,0,1,2\}\quad(i\ne j).
\]
The resulting invertibility domain after any subset of steps is
the initial \(I\) with the pairs \(\{y_j,y_j+1\}\) removed.
If, in addition,
\begin{equation}
 |y_i-y_j|>2,\qquad
 \sigma p_\rho(y_j)>0,\quad
 \sigma p_\rho(y_j+1)>0,\quad q(y_j)>0
                                                               \label{re-shifted-signs}
\end{equation}
for every \(\rho,j\), all the scalar weights in
\eqref{re-plus-scalar}--\eqref{re-minus-scalar} are positive.
\end{lemma}

\begin{proof}
Both signs remove exactly \(\{y_j,y_j+1\}\) from the
current \(I\).  The required points for a later plus step are
\(y_i,y_i+1\); for a later minus step they are
\(y_i,y_i+1,y_i+2\).  An intersection with a removed pair
would give one of the excluded integer differences.
The determinant multipliers from earlier steps, at a general
point \(Y\), are respectively
\[
 \frac{Y-y_j+1}{Y-y_j},\qquad
 \frac{Y-y_j-1}{Y-y_j}.
\]
At \(Y=y_i\) neither has a zero or a pole.  This proves
recursive admissibility.

For positivity, the two possible nontrivial multipliers of
\(x(Y)\) are
\[
 \frac{Y-y_j}{Y-y_j-1},\qquad
 \frac{Y-y_j-1}{Y-y_j}.
\]
When \(Y=y_i\) or \(y_i+1\), both are positive: the difference
\(Y-y_j\) is either greater than one or less than zero.
The two determinant multipliers at \(Y=y_i\) are positive
as well.  In both signs of the recursion, the required values
of \(x\) are exactly \(x(y_i)\) and \(x(y_i+1)\), and the
required value of \(q\) is \(q(y_i)\).
Thus every new scalar weight is a product or quotient of
positive numbers, proving the assertion inductively.
\end{proof}

This also gives the version with shifted spectral cuts used
in the source.  Namely, it is enough that for each \(j\)
either \(y_j\) or \(y_j+1\) satisfy
\eqref{rp-cut-margin}, and that the \(y_j\) are separated
by more than two.  Lemma~\ref{rp-parameter-signs} gives
\eqref{re-shifted-signs} in both cases, using \(q(Y-1)>0\)
when the cut is imposed at \(Y=y_j+1\).
For the additional domain point \(y_j+2\), its proof gives
slightly more than was stated: every factor \(Y+\rho_i\)
has absolute value at least \(5/2\) with a fixed sign,
so that sign persists for every shift of absolute value at
most two.  Consequently \(p_\rho(y_j+2)\ne0\).
No stronger separation between the unshifted parameters
\(a_j,b_j\) is needed for this version.

\begin{corollary}\label{re-density}
Let \(n\) be the number of steps.  If \(e\) is a choice
of \(0\) or \(1\) on each auxiliary factor, then
\begin{equation}
 x_n(Y)=\sum_{\rho,e}x_{\rho,e}(Y)P_\rho\otimes E_{ee},
 \qquad
 c_{{\rm wt},n}=\sum_{\rho,e}c_{\rho,e}P_\rho\otimes E_{ee},
 \qquad
 R_n=F_n^*w^n c_{{\rm wt},n}F_n .              \label{re-final-congruence}
\end{equation}
Initially \(x_{\rho,e}(Y)=\sigma p_\rho(Y)\) and
\(c_{\rho,e}=1\).
At a plus step of parameter \(a\), multiply
\(x_{\rho,e}(Y)\) by
\((Y-a)/(Y-a-1)\) when its new bit is \(0\), and by \(1\)
when it is \(1\).  Multiply \(c_{\rho,e}\) respectively by
\begin{equation}
                  q(a)/x_{\rho,e}(a+1),\qquad x_{\rho,e}(a).
                                                               \label{re-plus-scalar}
\end{equation}
At a minus step of parameter \(b\), the corresponding
multipliers are \((b-Y)/(b-Y-1)\), \(1\), and
\begin{equation}
                  q(b-1)/x_{\rho,e}(b-1),\qquad x_{\rho,e}(b).
                                                               \label{re-minus-scalar}
\end{equation}
Every quantity on the right is evaluated before the step.
Under the spectral-margin and separation hypotheses of
Corollary~\ref{rp-terminal-density-positive}, or under
the shifted hypotheses of Lemma~\ref{re-shifted-domain}, all
\(c_{\rho,e}>0\).
Consequently, for every smooth compactly supported radial
vector \(\Phi\),
\begin{equation}
 \int\langle R_n\Phi,\Phi\rangle
       =\int w^n\sum_{\rho,e}c_{\rho,e}
                       \|P_\rho(F_n\Phi)_e\|^2\geq0.             \label{re-positive-form}
\end{equation}
The integration measure is the radial measure specified at
the start of the section.
\end{corollary}

\begin{proof}
All initial projections commute, and every new formula for
\(x,c_{\rm wt}\) is diagonal on the new factor.
Induction therefore proves the spectral formulas and the
scalar recurrences.  They are exactly the scalar expressions
whose signs were proved in
Corollary~\ref{rp-terminal-density-positive}.
On a compact subset of \(G>0\), the inverses of \(G,\zeta\)
and \(\bar\zeta\), and all coefficients of the finitely many
operators, are bounded with all needed derivatives.
Differential operators do not enlarge supports.
Formal integration by parts is therefore legitimate in
\eqref{re-final-congruence}, giving
\eqref{re-positive-form}.  The pointwise orthogonality of
the \(P_\rho\) justifies its last expansion.
\end{proof}

For the empty sequence, \(n=0\), \(R_0=F_0=1\), and
\eqref{re-positive-form} is simply the ordinary squared norm.
When \(r=0\), there are no independent \(c_i\), and their
product measure and all assertions about their Euler
derivatives disappear.  When \(k=1\), the finite type has
one element, \(p_\rho=1\), and the initial correction
\(e(Y)\) is zero.  These cases require no additional
invertibility assumption.

This proves the actual algebraic frame and its compactly
supported positive form.  Passing to the specific noncompact
vectors used later still requires their regularity, decay,
support and operator-commutation hypotheses to be checked
against the analytic theorems of
Sections~\ref{jg-section} and \ref{jp-section}.
The algebraic congruence alone does not assert those hypotheses.
The correspondence of the formulas with the submitted source
is recorded separately; see \cite{alpo-uniform-source}.


